\section{Media project}

\subsection{Objectifs du media project}

L'objectif du media, est d'offrir une alternative entièrement textuelle aux moteurs de \gl{vn} facile à utiliser et à mettre en ligne.

En proposant une solution \gls{opensource}, nous laissons la possibilité aux personnes possédant des connaissances en développement de modifier et de rajouter des fonctionnalités à leurs projets. De plus, l'\gls{opensource} nous permet de laisser toutes personnes le souhaitant, de participer au développement du projet.



\subsection{Contraintes et limites}

Les moteurs de jeux, afin de rendre ces derniers jouables, compiles le code et le transforme en une version exécutable. Dans le cadre de ce média, nous allons réalisé un outil utilisant des technologies web, nous permettant de voir notre résultat sur n'importe quel navigateur internet. Nous expliquerons également comment mettre en ligne facilement et gratuitement les réalisations.

Une autre contrainte évoqué plus tôt, est de rendre le texte lisible et compréhensible autant pour les humains que pour les ordinateurs. Pour rendre cela possible, nous utiliserons le format \gls{json}. Ce format d'écriture clé/valeur, permet de stocker du texte attribué à des clés.

Pour des raisons de temps à attribué, pour les fonctionnalités nous nous contenterons de celles principales concernant les \gls{vn} qui sont les suivantes.


\begin{table}[htbp]
    \centering
    \caption{Fonctionnalités de bases du media project}
    \renewcommand{\arraystretch}{1.5}
    \begin{tabular}{p{4cm} p{9cm}}
        \hline
        \textbf{Fonctionnalité} & \textbf{Description} \\
        \hline
        Effet machine à écrire
        &
        Simule une machine à écrire, affichant le texte caractère après caractère.
        \\
        \addlinespace
        Gestion des scènes
        &
        Toutes les histoires sont divisées en scènes, le but ici est de permettre de naviguer d'une scène à une autre, d'effectuer un changement de décor.
        \\
        \addlinespace
        Gestion de saisies du joueur
        &
        Mettre un système permettant de stocké une saisie du joueur avec une valeur clé, facile à réutilisé.
        \\
        \addlinespace
        Gestion des choix
        &
        Permettre aux joueurs de choisir une réponse parmi plusieurs, les faisant changer de routes.
        \\
        \hline
    \end{tabular}
\end{table}

D'autres fonctionnalités pourraient être mises en place en fonction du temps, mais pour ce qui est des tests, nous allons nous focaliser sur ces fonctionnalités citées ci-dessus.